% Generated draft: requires mathematical and source review.
\documentclass[11pt]{article}
\usepackage[margin=25mm]{geometry}
\usepackage{fontspec}
\setmainfont{Latin Modern Roman}
\usepackage{amsmath,amssymb,mathtools,unicode-math}
\setmathfont{Latin Modern Math}
\usepackage{enumitem,needspace,xurl,hyperref}
\setlength{\parindent}{0pt}
\setlength{\parskip}{5pt}
\setlength{\emergencystretch}{3em}
\newcommand{\N}{\mathbb{N}}
\newcommand{\Z}{\mathbb{Z}}
\newcommand{\Q}{\mathbb{Q}}
\newcommand{\R}{\mathbb{R}}
\newcommand{\C}{\mathbb{C}}
\newcommand{\F}{\mathbb{F}}
\newcommand{\A}{\mathbb{A}}
\newcommand{\PP}{\mathbb{P}}
\newcommand{\E}{\mathbb{E}}
\newcommand{\eps}{\varepsilon}
\newcommand{\dd}{\,\mathrm{d}}
\begin{document}
\section*{UnsolvedMath ID: 800509}
\textbf{Source problem code:} catalogue entry
\subsection{Definitions and mathematical statement}
Let \(A\) denote the matrix to which a Lanczos algorithm is applied, and let \(q_k\) denote a vector used in its \(k\)-th matrix--vector multiplication. An inexact product can be written as
\[
\widehat y_k=Aq_k+e_k,
\]
where \(\widehat y_k\) is the returned product and \(e_k\) is its error. This notation alone imposes no restriction on the errors. The source does not specify the scalar field, matrix dimensions, symmetry or self-adjointness assumptions, or a particular variant of Lanczos.

\textbf{Research problem.} Analyze the Lanczos algorithm under general models of inexact matrix--vector multiplication, under the use of a preconditioner, or under both. Develop bounds that generalize Paige's analysis and/or provide new, easier-to-understand insight into the behavior of the algorithm. The source explicitly regards either type of contribution as progress.

\textbf{Unresolved specifications.} The source does not prescribe an admissible error model, its parameters or quantifiers, assumptions on a preconditioner, or how preconditioning enters the algorithm. It also does not identify the quantities to be bounded or the particular results of Paige to be generalized. These choices and the hypotheses needed for the selected Lanczos variant must be stated in any precise theorem addressing this research problem; no single fully quantified conjecture is supplied.
\subsection{Short English statement}
Study how inexact matrix--vector products and/or preconditioning affect Lanczos. Progress may consist of bounds extending Paige's analysis and/or clearer insight into the algorithm's behavior. The supplied record leaves the algorithmic variant, hypotheses, error model, preconditioning scheme, and target bounds unspecified.
\subsection{Sources}
\begin{itemize}
\item Local UnsolvedMath registry: permanent ID \(800509\), source problem code \texttt{catalogue entry}, title ``Linear Systems and Eigenvalue Problems --- Problem 18.'' The supplied record labels the problem open; this label is not an independently verified mathematical or literature assessment.
\item Supplied source attribution: Amsel, Baumann et al., \emph{Linear Systems and Eigenvalue Problems: Open Questions from a Simons Workshop} (2026), Problem 18. Supplied URL: \url{https://arxiv.org/abs/2602.05394}. The record reports access on 2026-10-16 and extraction from arXiv LaTeX; these provenance statements are supplied and have not been independently verified.
\item Paige's analysis is mentioned in the supplied statement, but no title, bibliographic reference, or specific theorem is supplied for it.
\end{itemize}
\end{document}
